\begin{proposition}[Conormal sections with trivial normalization pullback]\hypertarget{tag-01AZ}{}\label{cdp:st:A.5.1}
\StatementRecord{cdp:st:A.5.1}{proved}{01AZ}{\SeriesRef{cdp:st:A.1.1}, \SeriesRef{cdp:st:A.4.1}}

Let $S$ be a connected reduced full fibre of a proper map from a complex manifold to a smooth curve. For every line bundle $A$ on $S$ with $\eta^*A\simeq\OO_{\widetilde S}$, the group $H^0(S,V\otimes A)$ contains the canonical image of $H^0(\widetilde S,\eta^*(N^*\otimes A))$. A local base coordinate and a choice of trivialization of $\eta^*A$ identify this subspace with $H^0(\widetilde S,\OO_{\widetilde S})$. In particular it is nonzero. This applies to the subcase of \cite[Proposition~7.3, pp.~693--694]{CDP20} in which all $\eta^*(L_t|_S)$ are trivial. It is compatible with the vanishing of $H^0(S,Q\otimes A)$ through the maps in Proposition~\ref{cdp:st:A.1.1}, as \ref{cdp:st:A.4.1} demonstrates.
\end{proposition}


\begin{proof}
\begin{enumerate}[label={$\langle1\rangle$\arabic*.},leftmargin=*]
\item Proposition~\ref{cdp:st:A.2.2} gives the inclusion in
\eqref{cdp:eq:A.3}. Tensoring it by the invertible sheaf $A$ and applying
the finite projection formula gives
\[
 \eta_*\eta^*(N^*\otimes A)\hookrightarrow K\otimes A
                         \hookrightarrow V\otimes A.
\]
Locally, for a defining equation $F$ of $S$, the first map sends
$b\otimes[F]$ to $b\,dF|_S$. The inclusion of the Jacobian ideal in the
conductor makes its coefficients holomorphic on $S$. Under $F'=uF$,
both the conormal frame and $dF|_S$ are multiplied by $u$, so these maps
glue. Taking global sections gives the claimed canonical injection.

\item A coordinate $t$ about the image point on the base defines the full
fibre by $F=t\circ f$. Since this full fibre is reduced, $[F]$ is a frame
of $N^*$. Its pullback, together with the chosen trivialization of
$\eta^*A$, identifies the source with
$H^0(\widetilde S,\OO_{\widetilde S})$. The nonempty compact normalization
has the nonzero constant function $1$ on each connected component.
This proves nonvanishing and applies to each $A=L_t|_S$ in the stated
published subcase.

\item The image lies in $H^0(S,K\otimes A)$, the kernel of
$H^0(S,V\otimes A)\to H^0(S,Q\otimes A)$ by
\eqref{cdp:A:filtration2}. Thus these nonzero sections map to zero in
the quotient group. Proposition~\ref{cdp:st:A.4.1} gives the stated
example with vanishing quotient group.
\end{enumerate}
\end{proof}